\subsubsection{Datenschutz und Compilance}
\enquote{Data is the pollution problem of the information age, and protecting privacy is the environmental challenge.}
\parencite[S.~1]{schneier2015}
\newline
Dieses Zitat verdeutlicht die zunehmende Bedeutung des Themas Datenschutz in der heutigen Zeit, in der personenbezogene Daten in großem Umfang erhoben und verarbeitet werden. 
Auch im Rahmen des zu entwickelnden Bewerbermanagementportals werden sensible Daten verarbeitet, wie Namen und Adressen der Bewerber. 
Deshalb ist die Einhaltung datenschutzrechtlicher und regulatorischer Anforderungen eine zentrale Rolle bei der Systemkonzeption.
\newline
Das Bewerbermanagementsystem verarbeitet personenbezogene Daten wie Name, Kontaktdaten, Lebenslauf und weitere Bewerbungsunterlagen. 
Diese Daten duerfen ausschließlich zum Zweck der Durchführung des Bewerbungsverfahrens erhoben und verarbeitet werden. 
Dabei werden die Grundsätze der Zweckbindung, Datenminimierung und Transparenz berücksichtigt. 
Zudem wird sichergestellt, dass nur berechtigte Nutzerrollen (z. B. Recruiter:innen) Zugriff auf die Bewerberdaten erhalten. (vgl. Art.~5–6 \parencite{dsgvo}).
\newline
Vor dem Absenden einer Bewerbung müssen Bewerber:innen aktiv in die Verarbeitung ihrer personenbezogenen Daten einwilligen. 
Diese Einwilligung erfolgt über eine verpflichtende Checkbox im Bewerbungsformular, die auf die Datenschutzerklärung verweist. 
Dadurch wird sichergestellt, dass die Datenverarbeitung transparent erfolgt und auf einer klar dokumentierten Einwilligung der betroffenen Person basiert. (vgl. Art.~6 Abs. 1 lit a und 7 \parencite{dsgvo})
\newline
Bewerberdaten werden nur so lange gespeichert, wie es für den Bewerbungsprozess erforderlich ist. 
Nach Abschluss des Verfahrens werden personenbezogene Daten nach einer Frist automatisch gelöscht. 
Das es keine gestzlichen vorgegeben Fristen gibt, empfehlen Datenschutzbehoerden die Daten fuer maximal 6 Monate aufzubewahren. 
Somit sind die gesetzlichen Anforderungen zur Speicherbegrenzung und Datenminimierung eingehalten. (vgl Art.~5 Abs. 1 lit. \parencite{dsgvo})
\newline
Das Bewerbungsformular wird so gestaltet, dass keine diskriminierenden oder nicht erforderlichen personenbezogenen Informationen abgefragt werden. 
Angaben zu Merkmalen wie Geschlecht oder Alter werden nicht verpflichtend erhoben, um eine diskriminierungsfreie Bewerbungsprüfung zu gewährleisten. 
Damit wird sichergestellt, dass der Bewerbungsprozess den Grundsätzen der Gleichbehandlung entspricht. (§1 und §7 \parencite{agg})
